% =====================================================================
% CAPITULO 6 --- WRAPPER VISUAL
% Mantem o capitulo original e substitui apenas o banco complementar pela
% versao visual. Tambem repete a figura nas questoes que antes dependiam da
% expressao "mesmo circuito da questao anterior".
% =====================================================================
\begingroup

\def\fvSupCapDuasFontes{%
\begin{circuitikz}[american,scale=.86,transform shape,line width=.8pt]
  \coordinate (g) at (3.8,0);
  \coordinate (n) at (3.8,3.8);
  \draw (0,0) to[\fonteV,l^={$E_1=\SI{12}{\volt}$}] (0,3.8)
        to[R,l={$\SI{6}{\ohm}$}] (n);
  \draw (n) to[R,l_={$\SI{3}{\ohm}$}] (g) -- (0,0);
  \draw (7.6,0) to[\fonteV,l_={$E_2=\SI{6}{\volt}$}] (7.6,3.8)
        to[R,l_={$\SI{6}{\ohm}$}] (n);
  \draw (7.6,0)--(g);
  \node[Vcol,fill=white,inner sep=1.5pt] at ($(n)+(0,.58)$) {$V$};
  \node[ground] at(g){};
\end{circuitikz}}

\newcount\fvSupCapQuestao
\fvSupCapQuestao=0
\let\fvSupCapRespostaOriginal\resposta
\renewcommand{\resposta}[1]{%
  \global\advance\fvSupCapQuestao by 1\relax
  \ifnum\fvSupCapQuestao=4\relax\circuitoq{\fvSupCapDuasFontes}\fi
  \ifnum\fvSupCapQuestao=6\relax\circuitoq{\fvSupCapDuasFontes}\fi
  \fvSupCapRespostaOriginal{#1}}

\let\fvSupInputOriginal\input
\renewcommand{\input}[1]{%
  \ifstrequal{#1}{bancos/banco-teorema-da-superposicao.tex}%
    {\begingroup\let\input\fvSupInputOriginal
     \fvSupInputOriginal{bancos/banco-teorema-da-superposicao-visual.tex}%
     \endgroup}%
    {\fvSupInputOriginal{#1}}}
\fvSupInputOriginal{cap06-superposicao}
\endgroup
